\documentclass[../main.tex]{subfiles}
\begin{document}

Chapter 1 provides an overview of two key challenges in the field of fashion sketch-to-image translation: the limited availability of large-scale, well-annotated fashion sketch datasets and the high computational demands of existing state-of-the-art generative models. The chapter emphasizes the importance of constructing a comprehensive dataset that not only contains a substantial number of fashion sketches but also incorporates diverse sketch styles, levels of abstraction, and drawing characteristics to better reflect real-world design scenarios. It also reviews recent advances in sketch-to-image translation and motivates the development of lightweight yet effective generative approaches.
\section{Problem Statement}
\label{sec:problem_statement}
Fashion sketch-to-image generation has emerged as an important research topic due to its potential to bridge the gap between conceptual fashion design and realistic garment visualization \cite{cheng2021fashion}. In the conventional fashion design workflow, designers often begin with hand-drawn sketches to express their ideas before investing significant time and effort in digital rendering, prototyping, and sample production. The ability to automatically translate sketches into realistic clothing images can accelerate this process, provide rapid visual feedback, and support designers in exploring multiple design alternatives more efficiently. As a result, sketch-to-image generation has attracted increasing attention from both academia and industry as a promising tool for computer-aided fashion design \cite{isola2017image}.

Despite recent advances in generative models \cite{wang2018high}, developing an effective fashion sketch-to-image system remains a challenging task. Fashion sketches exhibit substantial variations in drawing style, abstraction level, line quality, and the amount of visual detail provided by different designers. These variations make it difficult for a model to accurately infer garment structures, textures, and precise color mappings from sparse sketch information. Furthermore, the availability of large-scale and diverse fashion sketch datasets remains limited \cite{sangkloy2016sketchy}, restricting the ability of existing models to generalize well across different sketch styles and real-world design scenarios.

Another important challenge lies in the computational complexity of modern image generation models \cite{menghani2023efficient}. While recent state-of-the-art approaches have achieved remarkable visual quality, many of them require substantial computational resources, long training times, and expensive hardware for inference. Such requirements limit their applicability in practical fashion design environments, where fast response time and computational efficiency are often as important as image quality. This limitation is particularly relevant for interactive design systems that aim to provide near real-time feedback during the creative process.

Given these challenges, there is a need for sketch-to-image generation approaches that can effectively handle diverse fashion sketch inputs while maintaining a favorable balance between visual quality and computational efficiency. Therefore, this thesis focuses on the fashion sketch-to-image translation problem and investigates lightweight generative models capable of producing realistic garment images from sketches with reduced computational cost. The proposed research aims to contribute toward practical and accessible computer-aided fashion design systems that can support rapid design iteration and improve workflow efficiency for both professional and novice designers.

\section{Background and Problems of Research} 
\label{sec:background}

\subsection{Review of Existing Approaches}
To address the fashion sketch-to-image generation problem outlined in Section \ref{sec:problem_statement}, numerous approaches have been proposed, primarily leveraging image-to-image translation frameworks. Foundational models such as Pix2Pix \cite{isola2017image} established the viability of conditional Generative Adversarial Networks (cGANs) for translating edge maps into natural images by relying on strictly paired datasets. To overcome the dependency on paired training data, models like CycleGAN \cite{zhu2017unpaired} were introduced, utilizing cycle-consistency loss to learn domain mappings without requiring one-to-one image correspondences. Moreover, in order to enhance generation diversity, Multimodal Unsupervised Image-to-Image Translation (MUNIT) \cite{huang2018multimodal} was developed, allowing the generation of diverse outputs from a single sketch by disentangling content and style representations.

Further advancing the field, models based on StyleGAN \cite{karras2019style} inversion, such as pixel2style2pixel (pSp) \cite{richardson2021encoding}, have been introduced to directly map sketch inputs into a pre-trained StyleGAN latent space, enabling the generation of highly realistic and controllable images without requiring task-specific architectures. More recently, diffusion models such as Stable Diffusion \cite{rombach2022high} have emerged as a powerful alternative for conditional image synthesis, offering unprecedented generation quality and diversity.

\subsection{Limitations of Current Solutions}
While recent studies have demonstrated impressive visual realism, several critical limitations restrict their applicability in practical fashion design systems. 

First, the pursuit of higher visual fidelity has led to increasingly complex architectures. State-of-the-art diffusion-based models often comprise hundreds of millions or even billions of parameters, resulting in substantial inference latency and memory consumption. Such computational demands are often at odds with the requirements of interactive design environments, where near real-time feedback and deployment on resource-constrained devices are essential. Comparable challenges can also be observed in GAN-based inversion frameworks. For example, while pSp can generate highly realistic outputs, it relies on a computationally intensive StyleGAN generator and a complex inversion encoder, leading to increased inference time.

Second, existing methodologies often struggle to balance structural preservation with diverse generation. Unpaired models like CycleGAN frequently fail to retain the fine-grained geometric details specified in a fashion sketch. Multimodal frameworks like MUNIT address diversity but often struggle to cleanly disentangle structural content from color variations, leading to style artifacts or structural inconsistencies.

Furthermore, current solutions are typically constrained by the lack of optimized datasets; publicly available sketch datasets are often limited in scale or lack the detailed color annotations necessary to effectively train and evaluate robust, color-conditioned generative frameworks.

\section{Research Objectives and Conceptual Framework}
\label{sec:objectives}
To address the aforementioned challenges, this thesis outlines two primary objectives: first, to construct a comprehensive, annotated dataset tailored for fashion sketch-to-image tasks and second, to develop a lightweight, highly controllable generative model that balances visual fidelity, structural preservation, and computational efficiency. 

To realize these objectives, this thesis proposes a conceptual framework that provides specific, targeted solutions to overcome the three primary limitations discussed in Section \ref{sec:background}:

\begin{enumerate}
    \item Overcoming computational complexity: To tackle the latency and memory bottlenecks inherent in large-scale diffusion models and complex GAN inversions, the framework introduces the MS2I (Mobile Sketch-to-Image) architecture. This novel model incorporates Reparameterizable Convolutions and Monarch Attention mechanisms to significantly reduce the parameter count and computational overhead during inference, ensuring the model remains lightweight and suitable for deployment in resource-constrained environments requiring rapid visual feedback.
    
    \item Balancing structural preservation with generation diversity: Further building upon the MS2I architecture to rectify structural inconsistencies and style entanglement observed in prior unpaired and multimodal networks, the model employs a graduated style modulation approach. This strategy allows MS2I to faithfully preserve the fine geometric details specified in the input sketch while cleanly disentangling structural features from color representations, enabling diverse and highly controllable image synthesis.
    
    \item Addressing the lack of suitable datasets: To resolve the data scarcity issue, the framework focuses on constructing a large-scale, high-quality paired dataset. By leveraging publicly available fashion image repositories, particularly Fashionpedia, and employing advanced edge-based sketch generation techniques, a comprehensive dataset is synthesized. To better reflect real-world fashion sketching practices, the generated sketches incorporate diverse drawing styles and line characteristics, ranging from clean contour-based sketches to more abstract representations. This dataset is further enriched with comprehensive color annotations, providing a robust foundation for training and evaluating conditional generative models.
\end{enumerate}

\section{Contributions}
\label{sec:contributions}
The main contributions of this thesis are as follows:
\begin{enumerate}
    \item The construction of a large-scale, high-quality dataset comprising approximately 66,000 paired sketch-fashion image samples, enriched with diverse sketching styles and detailed color annotations.
    \item The development and implementation of a lightweight GAN-based model capable of generating realistic fashion images from sketches with reduced computational overhead and latency.
\end{enumerate}

\section{Organization of Thesis}
\label{sec:organization}
The remainder of this thesis is organized as follows:

Chapter 2 discusses the scope of the thesis and  provides a comprehensive literature review on existing sketch-to-image generation architectures and related fashion datasets. It examines the fundamental background knowledge, evaluating the strengths and critical limitations of current approaches to motivate the proposed solutions.

Chapter 3 details the proposed methodology and the core technical contributions. It first describes the data preparation pipeline used to construct the large-scale, paired fashion sketch dataset with detailed color annotations. Second, it introduces the overall architectural design of the proposed model, outlining its lightweight structural components and the mechanisms employed for effective style modulation.

Chapter 4 presents a theoretical analysis of the proposed architecture. It focuses on the model's computational complexity and analyzes the underlying mechanisms—specifically how the parameter reduction techniques enable lightweight and highly efficient image generation suitable for resource-constrained environments.

Chapter 5 reports the quantitative and qualitative results from extensive experiments. It outlines the experimental setup and evaluation metrics, followed by a comprehensive comparative analysis of the proposed model against prominent GAN-based models. This chapter demonstrates the model's capability to achieve high-quality, color-controlled generation with significantly reduced latency.

Chapter 6 concludes the thesis by summarizing the main findings and technical contributions, and discussing practical implications for interactive fashion design. It also discusses the remaining challenges and limitations encountered throughout the research before outlining potential directions for future work.

\end{document}